% =========================================================================
% PROVENANCE (section-level)
%   status      : CODE_FACT unless marked otherwise
%   produced-by : utilities.py :: classify_mc_dict()  (lines 519-635)
%                 utilities.py :: DecayMode_new       (line 73)
%                 utilities.py :: create_templates_new (lines 1073-1273)
%   commit      : dd520a5
%   scope       : MC16rd, Run 1 + Run 2, e and mu channels
% =========================================================================
\section{Truth Classification and Simulated Sample Composition}
\label{sec:truth}

Every simulated candidate is assigned to exactly one truth category by
\texttt{utilities.classify\_mc\_dict()}.
The categories are the samples from which the fit templates of
Sec.~\ref{sec:fit} are built, so their definitions, their ordering, and their
completeness are part of the statistical model rather than bookkeeping.

\subsection{Classification variables}
\label{sec:truth:vars}

The decision tree uses four truth quantities:
\texttt{D\_mcErrors}, the MC-matching error flags of the $D$ candidate;
the lepton MC PDG code (\texttt{ell\_BFbrems\_mcPDG} in the electron channel,
i.e.\ the \emph{pre}-bremsstrahlung-correction match, and
\texttt{ell\_mcPDG} in the muon channel);
\texttt{B0\_isContinuumEvent}; and the MC PDG codes of the $B^{0}$ candidate
and of its first two generated daughters.

Two derived predicates organise the tree:
\begin{align}
  \mathrm{trueD} &\equiv \texttt{D\_mcErrors} = 0, \\
  \mathrm{fakeD} &\equiv 0 < \texttt{D\_mcErrors} < 512, \\
  \mathrm{fakeTracks} &\equiv \texttt{D\_mcErrors} = 512 ,
\end{align}
where $512$ is the clone/fake-track bit.
True and fake leptons are separated by whether the matched lepton PDG code
has the absolute value expected for the channel.

\subsection{Decision tree and precedence}
\label{sec:truth:tree}

The categories are assigned in the order below.
Fake-object categories take precedence over every category built from a true
$D$ and a true lepton, and, importantly, they are assigned \emph{without} a
requirement on \texttt{B0\_isContinuumEvent}: a continuum event containing a
fake $D$ is classified as \texttt{bkg\_fakeD}, not as \texttt{bkg\_continuum}.
The continuum flag is used only to separate the categories that already have
a true $D$ and a true lepton.

\begin{enumerate}
  \item \texttt{bkg\_fakeD} --- $0 < \texttt{D\_mcErrors} < 512$.
        The $K\pi\pi$ combination does not come from a real $D^{\pm}$.
  \item \texttt{bkg\_fakeL} --- trueD and the lepton is not a true lepton of
        the channel flavour.
  \item \texttt{bkg\_fakeTracks} --- $\texttt{D\_mcErrors} = 512$: one of the
        $D$ daughters is a clone or fake track.
\end{enumerate}

The remaining candidates have a true $D$ and a true lepton (denoted TDT$\ell$)
and are split by the origin of the pair:

\begin{enumerate}\setcounter{enumi}{3}
  \item \texttt{bkg\_continuum} --- $\texttt{B0\_isContinuumEvent} = 1$:
        the event is $e^{+}e^{-} \to q\bar{q}$ rather than
        $\Upsilon(4S) \to B\bar{B}$.
  \item \texttt{bkg\_combinatorial} --- $\texttt{B0\_mcPDG} = 300553$:
        the $D$ and the lepton are both real but come from \emph{different}
        $B$ mesons of the same $B\bar{B}$ event.
  \item \texttt{bkg\_hadronicB\_secondaryL} --- a single, non-continuum $B$
        decay in which the lepton is \emph{not} a direct daughter of the
        reconstructed $B$ and does not come from a $\tau$ of that $B$.
        This is the double-charm/secondary-lepton family: hadronic decays
        such as $B \to D D_{(s)} X$ in which one charm meson decays
        semileptonically and supplies the lepton.
  \item Signal-like decays (\texttt{signals}) --- $|\texttt{B0\_mcPDG}| \in
        \{511, 521\}$ \emph{and} the lepton is either a direct daughter of
        that $B$ or the daughter of a $\tau$ from that $B$.
  \item \texttt{bkg\_other\_TDT$\ell$} --- explicit catch-all: any TDT$\ell$
        candidate matched by none of items 4--7.
\end{enumerate}

The signal-like sample is then subdivided by the generated first two
daughters of the $B$:

\begin{enumerate}\setcounter{enumi}{8}
  \item $B \to D \tau \nu$ (signal) and $B \to D \ell \nu$ (normalisation),
        identified by the PDG-code product of the first two daughters;
  \item $B \to D^{*} \tau \nu$ and $B \to D^{*} \ell \nu$;
  \item $B \to D^{**} \tau \nu$, and $B \to D^{**} \ell \nu$ split into
        narrow ($D_{1}, D_{2}^{*}$) and broad ($D_{0}^{*}, D_{1}'$) states;
  \item semileptonic ``gap'' modes, identified by a charged $D^{(*)}$
        together with a $\pi$ or an $\eta$, merged into a single
        \texttt{$D\ell\nu$\_gap} category;
  \item \texttt{bkg\_other\_signal} --- explicit catch-all: any signal-like
        candidate matched by none of items 9--12.
\end{enumerate}

Each category is finally stamped with an integer code from
\texttt{DecayMode\_new}, which is the interface used by the downstream
template and workspace code.

\subsection{Are the categories mutually exclusive and exhaustive?}
\label{sec:truth:closure}

\paragraph{Exclusivity.}
The three fake categories are disjoint by construction, because they
partition disjoint ranges of \texttt{D\_mcErrors} and the fake-lepton
category additionally requires $\texttt{D\_mcErrors} = 0$.
Among the TDT$\ell$ categories, exclusivity is \emph{inherited from an
invariant of the MC truth record} rather than enforced by the queries
themselves: \texttt{bkg\_continuum} places no requirement on
\texttt{B0\_mcPDG}, while \texttt{bkg\_combinatorial} and the signal-like
selection place no requirement on \texttt{B0\_isContinuumEvent}.
These definitions are disjoint precisely because
$\texttt{B0\_isContinuumEvent} = 1$ is a necessary and sufficient condition
for a continuum event, so that \texttt{B0\_mcPDG} is meaningful only when the
flag is zero, and a candidate with $|\texttt{B0\_mcPDG}| \in \{511,521\}$ or
$\texttt{B0\_mcPDG} = 300553$ necessarily has the flag equal to zero
(\texttt{ANALYSIS\_DECISION}).
The categories are therefore exclusive, but the guarantee is a property of
the simulation, not of the code; a change in the meaning of either variable
would silently break it.

\paragraph{Exhaustiveness.}
The TDT$\ell$ and signal-like branches are exhaustive by construction,
because each terminates in an explicit catch-all
(\texttt{bkg\_other\_TDT$\ell$} and \texttt{bkg\_other\_signal}) built as the
complement of the categories above it.
The fake branch has no such catch-all.
Candidates with $\texttt{D\_mcErrors} > 512$ --- that is, the clone/fake-track
bit set together with any other mismatch bit, or any higher bit --- satisfy
none of trueD, fakeD or fakeTracks, and are therefore assigned to no category
at all.
Such candidates are expected to be absent or negligible in MC16rd
\AnalysisTBD{to be confirmed: run
\texttt{scripts/validate\_truth\_categories.py} over one MC16rd file and
report the number of unclassified candidates and their
\texttt{D\_mcErrors} spectrum.  Until then, exhaustiveness of the fake branch
is an assumption}

\paragraph{Categories excluded from the fit.}
Three classified categories are dropped before the templates are built
(\texttt{create\_templates\_new}, default
\texttt{sample\_to\_exclude}): \texttt{bkg\_fakeTracks},
\texttt{bkg\_other\_TDT$\ell$} and \texttt{bkg\_other\_signal}.
The classification is therefore exhaustive in a stronger sense than the fit
model is: the two catch-alls exist to make the classification complete, and
are then excluded from the likelihood.
\AnalysisTBD{quantify the yield of the three excluded categories in the
signal region and in the $D$-mass sideband, and either justify neglecting
them or assign a systematic uncertainty for doing so}

\subsection{Category merging in the fit templates}
\label{sec:truth:merging}

Two merges happen between classification and template construction, so the
category list of Sec.~\ref{sec:truth:tree} is not identical to the sample
list of Sec.~\ref{sec:fit}:
the narrow and broad $D^{**}\ell\nu$ categories are summed into a single
$D^{**}\ell\nu$ template; and when the gap-mode sample weight is set to zero
the combined template is relabelled $D^{**}\ell\nu + \text{gap}$, which folds
the gap modes into the $D^{**}$ template rather than dropping them.
\AnalysisTBD{state which of these configurations is nominal, and what
information is lost by merging narrow and broad $D^{**}$ given that they have
different form factors and different $\mmiss$ shapes}

\subsection{Generic-$B\bar{B}$ subdivision}
\label{sec:truth:bbbar}

The two generic-$B\bar{B}$ categories --- \texttt{bkg\_combinatorial} and
\texttt{bkg\_hadronicB\_secondaryL} --- are further subdivided for the
purpose of the reweighting described in Sec.~\ref{sec:bbbar}.
That subdivision is a property of the reweighting, not of the truth
classification, and is documented there.

\subsection{Sample composition}
\label{sec:truth:composition}

\AnalysisTBD{table of expected yields and fractions per category, in the
signal region and in the $D$-mass sideband, separately for Run 1 and Run 2
and for the $e$ and $\mu$ channels, at the nominal working point.  Producing
artifact and commit to be recorded in the provenance block above the table}
